\rhead{RB-EL1: Advanced Control Systems}
\phantomsection 
\section*{Advanced Control Systems (RB-EL1)}
\label{elec:RB_1}

\noindent \small{\hyperref[main:scheme]{$\leftarrow$ Back to Scheme}} \vspace{0.5cm}

\noindent \textbf{Credits:} 4 (3-1-0)

\subsection*{Theory}
\noindent\textbf{Unit 1} \\
\textit{State-Space Analysis and Design:} State-space representation (controllable canonical, observable canonical, diagonal forms), Controllability and observability tests (Popov-Belevitch-Hautus eigenvalues), Pole placement by state feedback, Ackermann's formula, State observers (Luenberger full-order/reduced-order), Separation principle, Riccati equation fundamentals.

\vspace{0.3cm}

\noindent\textbf{Unit 2} \\
\textit{Optimal Control Systems:} Linear Quadratic Regulator (LQR) design (finite/infinite horizon), Algebraic Riccati Equation (ARE) solution, Steady-state LQR and optimal gain computation, Linear Quadratic Gaussian (LQG) control, Kalman filtering for state estimation, Loop transfer recovery (LQR/LQG trade-offs), Disturbance decoupling and regulation.

\vspace{0.3cm}

\noindent\textbf{Unit 3} \\
\textit{Robust Control Design:} H-infinity control fundamentals (small gain theorem, bounded real lemma), Mixed-sensitivity optimization, mu-synthesis for structured uncertainty, Loop-shaping design procedure, Robust stability/performance margins, Kharitonov theorem for interval plants, Quantitative feedback theory (QFT) basics.

\vspace{0.3cm}

\noindent\textbf{Unit 4} \\
\textit{Adaptive Control Systems:} Model Reference Adaptive Control (MRAC: MIT rule, Lyapunov-based), Self-tuning regulators (explicit/implicit), Parameter convergence analysis, Persistent excitation conditions, Dead-zone modification for noise robustness, Gain scheduling strategies, Switching and supervisory control, Adaptive pole placement.

\vspace{0.3cm}

\noindent\textbf{Unit 5} \\
\textit{Nonlinear and Modern Control:} Feedback linearization (input-state, input-output), Sliding mode control (reaching law, boundary layer), Backstepping design for strict-feedback systems, Lyapunov stability analysis (direct/indirect methods), Passivity-based control, Model Predictive Control (MPC) fundamentals, Constrained optimization (QP solvers).

\newpage

\subsection*{Practical/Simulations}
\noindent\textbf{Session 1: State-Space Modeling} \\
Focuses on MATLAB state-space conversion, controllability/observability analysis.

\vspace{0.2cm}

\noindent\textbf{Session 2: LQR/LQG Controller Design} \\
Focuses on Riccati equation solution, Kalman filter tuning, simulation validation.

\vspace{0.2cm}

\noindent\textbf{Session 3: H-Infinity Synthesis} \\
Focuses on hinfsyn toolbox, mixed-sensitivity design, robust stability margins.

\vspace{0.2cm}

\noindent\textbf{Session 4: Adaptive Control Simulation} \\
Focuses on MRAC implementation, parameter tracking, robustness testing.

\vspace{0.2cm}

\noindent\textbf{Session 5: Nonlinear Control Design} \\
Focuses on feedback linearization, sliding mode simulation, Lyapunov validation.

\vspace{0.2cm}

\noindent\textbf{Session 6: MPC Implementation} \\
Focuses on quadprog solver, receding horizon control, constraint handling.

\vspace{0.2cm}

\noindent\textbf{Session 7: Robust Control Applications} \\
Focuses on QFT toolbox, uncertainty modeling, performance specification.

\vspace{0.2cm}

\noindent\textbf{Session 8: Real-Time Control Experiments} \\
Focuses on dSPACE/Simulink Real-Time, hardware-in-loop validation.

\vspace{0.2cm}

\noindent\textbf{Session 9: System Identification Integration} \\
Focuses on subspace methods feeding control design.

\vspace{0.2cm}

\noindent\textbf{Session 10: Control System Performance Analysis} \\
Focuses on comparative benchmarking, sensitivity functions, robustness metrics.

\vspace{0.2cm}

\newpage
\rhead{Curriculum Schema}
